\section*{Sets and Indices}

\begin{itemize}
    \item $C$: set of candidates (from \texttt{table\_1.csv}); here
    \[
    C=\{F,G,H,I,J\}.
    \]
    \item $c \in C$: candidate index.
    \item $P \subseteq C$: special conflicting candidate subset for the pairwise hiring restriction; here
    \[
    P=\{G,J\}.
    \]
\end{itemize}

\section*{Parameters}

\begin{itemize}
    \item $s_c$: salary of candidate $c$.
    \item $k_c$: skill level of candidate $c$.
    \item $e_c$: project management experience of candidate $c$.
    \item $B$ (code data: \texttt{budget}): company hiring budget.
    \item $M$ (code data: \texttt{max\_employees}): maximum number of employees that can be hired.
    \item $K^{\min}$ (code data: \texttt{min\_skill}): minimum required total skill level.
    \item $E^{\min}$ (code data: \texttt{min\_pm\_exp}): minimum required total project management experience.
    \item $U_{GJ}$ (code data: \texttt{max\_gj}): maximum number of hires allowed among candidates $G$ and $J$.
\end{itemize}

For this instance, the candidate-specific values are:
\[
\begin{array}{c|ccccc}
c & F & G & H & I & J \\ \hline
s_c & 12000 & 15000 & 18000 & 5000 & 10000 \\
k_c & 2 & 3 & 4 & 1 & 2 \\
e_c & 1 & 2 & 2 & 5 & 4
\end{array}
\]
and
\[
B=40000,\quad M=4,\quad K^{\min}=8,\quad E^{\min}=8,\quad U_{GJ}=1.
\]

\section*{Decision Variables}

\begin{itemize}
    \item $x_c$ (code: \texttt{x[c]}) for each $c \in C$, where
    \[
    x_c=
    \begin{cases}
    1 & \text{if candidate } c \text{ is hired},\\
    0 & \text{otherwise.}
    \end{cases}
    \]
\end{itemize}

\section*{Objective}

Minimize total salary of the hired team:
\[
\min \sum_{c \in C} s_c x_c.
\]

\section*{Constraints}

\begin{align}
\sum_{c \in C} s_c x_c &\le B
&& \text{(budget)} \\
\sum_{c \in C} x_c &\le M
&& \text{(maximum number of hires)} \\
\sum_{c \in C} k_c x_c &\ge K^{\min}
&& \text{(minimum total skill)} \\
\sum_{c \in C} e_c x_c &\ge E^{\min}
&& \text{(minimum total project management experience)} \\
x_G + x_J &\le U_{GJ}
&& \text{(at most one of $G$ and $J$)} \\
x_c &\in \{0,1\}
&& \forall c \in C.
\end{align}

For this instance, the explicit model is:
\begin{align}
\min\quad &12000x_F+15000x_G+18000x_H+5000x_I+10000x_J \\
\text{s.t.}\quad
&12000x_F+15000x_G+18000x_H+5000x_I+10000x_J \le 40000 \\
&x_F+x_G+x_H+x_I+x_J \le 4 \\
&2x_F+3x_G+4x_H+x_I+2x_J \ge 8 \\
&x_F+2x_G+2x_H+5x_I+4x_J \ge 8 \\
&x_G+x_J \le 1 \\
&x_F,x_G,x_H,x_I,x_J \in \{0,1\}.
\end{align}

\section*{Notes on Implementation}

\begin{itemize}
    \item The implementation in \texttt{current\_heuristic.py} builds exactly the above binary linear optimization model using one binary variable per candidate.
    \item The objective and the budget constraint use the same salary expression; this is intentional and should be preserved, since the code includes both.
    \item The special incompatibility restriction is implemented only for candidates \texttt{G} and \texttt{J} via \texttt{x['G'] + x['J'] <= max\_gj}; no other mutual-exclusion structure is modeled.
    \item The model is solved by a MIP solver through \texttt{prob.solve(GUROBI\_CMD(msg=0))}, i.e., as a standard binary linear program rather than by custom enumeration or dynamic programming.
    \item Although all parameter values are read as numeric data from CSV files, \texttt{max\_employees} and \texttt{max\_one\_candidate\_g\_j} are explicitly cast to integers in code before being used.
\end{itemize}
